\chapter{Framework Methodology}
\label{ch:framework}
\label{ch:edge_computing}
\label{ch:federated_learning}
\label{ch:communication}
\label{ch:root_cause_analysis}
\label{ch:adaptive_thresholding}

This chapter presents the methodological core of the proposed framework. The
design combines edge-local anomaly detection, federated model training,
differential privacy, compressed communication, causal root cause analysis, and
adaptive threshold control. These mechanisms are presented together because they
operate as one feedback loop: local agents detect anomalies, the coordinator
aggregates model updates, RCA contextualizes alerts, and threshold policies are
updated from operational feedback.

\section{Edge Telemetry Pipeline}
\label{sec:metrics_collection}

Each monitored service is paired with an edge agent. The agent collects
resource, network, and application metrics such as CPU utilization, memory
pressure, request latency, error rate, queue depth, and connection count. These
features are chosen because infrastructure-only metrics often miss
service-level degradation, while application-only metrics may not reveal the
resource cause behind a symptom.

The set of collected metrics is not inferred automatically by the model; it is
defined by the observability contract of the deployment. In a production
microservice platform, resource metrics can be exported by Prometheus
node/process exporters, application metrics by service instrumentation or
StatsD-compatible receivers, and call relationships by OpenTelemetry traces. In
the reproducible experiments, the SMD benchmark provides 38 numerical server
signals rather than named Prometheus metrics. The thesis therefore uses the full
38-dimensional SMD input and interprets the channels through the same feature
roles used by the framework: saturation, network pressure, runtime behavior, and
service health. This mapping keeps the benchmark reproducible while avoiding the
claim that every production deployment exposes exactly the same metric names.

Raw metrics are normalized locally using rolling statistics and segmented into
overlapping windows. For feature $j$ at time $t$, z-score normalization is:

\begin{equation}
    x'_{tj} = \frac{x_{tj} - \mu_{j,t}}{\sigma_{j,t} + \epsilon},
    \label{eq:local_normalization}
\end{equation}

where $\mu_{j,t}$ and $\sigma_{j,t}$ are computed over a recent local history
and $\epsilon$ prevents division by zero. For reproduction, $x_{tj}$ denotes
the raw value of feature $j$ at sample time $t$, $\mu_{j,t}$ and
$\sigma_{j,t}$ are the mean and standard deviation over the local rolling
history before the current decision, and $x'_{tj}$ is the normalized value
stored in the model input buffer. A window
$X_t \in \mathbb{R}^{W \times F}$ is then formed from the last $W=50$
observations and $F=38$ SMD channels in the evaluation. With the default 10\,s
scrape interval, one input window represents approximately 8.3 minutes of recent
behavior. Keeping this preprocessing at the edge has three benefits: local
inference can continue during coordinator outages, raw telemetry does not leave
the node, and the network carries only model updates and anomaly summaries.

The benchmark channels cover four operational roles: resource saturation (CPU,
memory, disk), network pressure (throughput, connections, retransmission),
application health (latency, error rate, queue depth), and runtime internals
(thread count, pool pressure). This is a feature contract: deployments may use
different metric names, but each feature should play one of these roles.

Preprocessing is conservative. Short missing intervals are forward-filled;
longer gaps are masked. Rolling statistics are stored per service so that the
same absolute latency can be treated differently for different services. This
per-service normalization also prevents leaking raw scale information through
centrally computed preprocessing constants in the federated setting.

\section{LSTM-Autoencoder Detector}
\label{sec:method_lstm_ae}

The anomaly detector is an LSTM-Autoencoder trained to reconstruct normal
telemetry windows. The encoder maps the input sequence to a latent state:

\begin{equation}
    z_t = Enc_\theta(X_t),
    \label{eq:encoder}
\end{equation}

and the decoder reconstructs the original sequence:

\begin{equation}
    \hat{X}_t = Dec_\phi(z_t).
    \label{eq:decoder}
\end{equation}

The anomaly score is the mean reconstruction error:

\begin{equation}
    e_t = \frac{1}{WF}\sum_{i=1}^{W}\sum_{j=1}^{F}
    \left(X_{t,ij} - \hat{X}_{t,ij}\right)^2 .
    \label{eq:reconstruction_error}
\end{equation}

The LSTM gates (Eqs.~\ref{eq:forget_gate}--\ref{eq:hidden_state} in
Chapter~\ref{ch:model_architecture}) regulate which temporal patterns are
retained. This is important for operational telemetry because anomalies can
depend on several earlier observations: a memory leak is a slow drift, while
a connection storm is a rapid multi-feature transition.

Training is unsupervised and minimizes reconstruction loss on windows treated as
normal. At inference time, a window is anomalous when:

\begin{equation}
    y_t =
    \begin{cases}
    1, & e_t > \tau_t,\\
    0, & e_t \leq \tau_t,
    \end{cases}
    \label{eq:anomaly_decision}
\end{equation}

where $\tau_t$ is a per-service threshold. The initial threshold is set from the
95th percentile of local training errors, then passed to the adaptive
thresholding component.
The architecture intentionally avoids a bidirectional LSTM in the edge detector.
Bidirectional recurrence can improve offline sequence modeling, but it requires
the full window before producing a representation and roughly doubles recurrent
computation. For streaming monitoring, the unidirectional LSTM-AE offers a
better latency profile and aligns with the causality of online telemetry: the
detector should decide from past and present observations, not from future
samples that would not yet be available in production.

\section{Federated Learning Protocol}
\label{sec:fl_architecture}

The framework uses FedAvg (Eq.~\ref{eq:bg_fedavg}) to train a global LSTM-AE
without centralizing raw telemetry. At the beginning of round $t$, the
coordinator exposes $\theta_t$. Participating edge agents download it, train
locally for $E=5$ epochs, and submit a parameter update weighted by local
dataset size $n_k$. Partial participation is supported: the coordinator
aggregates updates arriving before the round timeout, provided the minimum
threshold is met.

Three design choices suit unreliable edge deployments: (1)~pull-based model
rollout, so agents request the latest checkpoint rather than depending on a
broadcast; (2)~write-ahead round logging, so restarts never expose a partially
written model; and (3)~version-bounded staleness, so updates from outdated
models are discarded.

\begin{table}[H]
    \centering
    \small
    \begin{tabular}{|p{0.23\textwidth}|p{0.29\textwidth}|p{0.36\textwidth}|}
        \hline
        \textbf{Stage} & \textbf{Failure mode} & \textbf{Mitigation} \\
        \hline
        Model download & coordinator unavailable & continue inference with
        cached model \\
        \hline
        Local training & node crash & replay pending update in standard
        profile; restart in lightweight profile \\
        \hline
        Upload & timeout or network partition & join the next round; do not
        block aggregation \\
        \hline
        Aggregation & insufficient participants & abort round if
        $K_{\min}$ is not met \\
        \hline
    \end{tabular}
    \caption{Fault-aware behavior in the federated training protocol}
    \label{tab:fl_failures}
\end{table}

\section{Differential Privacy Mechanism}
\label{sec:method_dp}

Differential privacy is applied before model updates leave the edge node. Each
per-sample gradient $g_i$ is clipped to a maximum norm $C$:

\begin{equation}
    \bar{g}_i =
    g_i \cdot \min\left(1,\frac{C}{\|g_i\|_2}\right).
    \label{eq:method_clip}
\end{equation}

The batch update is then perturbed with Gaussian noise:

\begin{equation}
    \tilde{g} =
    \frac{1}{B}\sum_{i=1}^{B}\bar{g}_i
    + \mathcal{N}(0,\sigma^2 C^2 I).
    \label{eq:method_noise}
\end{equation}

In Eq.~\ref{eq:method_clip}, $g_i$ is the gradient contribution of training
sample $i$, $\|g_i\|_2$ is its Euclidean norm, and $C$ is the clipping bound.
If the norm is already below $C$, the gradient is unchanged; otherwise it is
rescaled to lie on the radius-$C$ ball. In Eq.~\ref{eq:method_noise}, $B$ is
the local batch size, $\tilde{g}$ is the noisy gradient used by the optimizer,
$\sigma$ is the DP noise multiplier, and $I$ is the identity covariance matrix,
so independent Gaussian noise is added to each gradient coordinate. These
definitions are necessary for reproduction because using batch-level clipping
instead of per-sample clipping, or adding noise before averaging instead of
after clipping, changes both utility and privacy accounting.

The evaluation (\S\ref{subsec:eval_dp_impact}) uses this mechanism to benchmark
the privacy--utility frontier. The key finding is that clipping is a decisive
control variable: reducing $C$ shrinks the effective noise $\sigma C$, recovering
utility but also reducing the practical perturbation. The thesis therefore
treats clipping as both a privacy and utility lever, and reports the residual
privacy concern alongside the utility result.

\section{Communication Reduction}
\label{sec:compression}

Federated training exchanges model states, so communication efficiency depends
on serialization and compression. The framework distinguishes three layers:

\begin{itemize}
    \item \textbf{serialization}: Torch binary serialization is used for Python
    model checkpoints, while Protocol Buffers define the typed gRPC control
    messages;
    \item \textbf{algorithmic compression}: optional Top-K sparsification keeps
    the largest-magnitude update components and accumulates the residual
    locally for later rounds \cite{aji2017sparse, karimireddy2019errorfeedback};
    \item \textbf{byte-level compression}: Zstandard is used when bandwidth is
    the bottleneck, while LZ4 is used when CPU headroom is low
    \cite{collet2018zstandard, lz4_reference}.
\end{itemize}

The adaptive compression policy is intentionally simple. If edge CPU utilization
is below 80\%, the agent selects Zstd to reduce bytes on the wire. Above that
threshold, the agent selects LZ4 to minimize compression latency. The evaluation
shows that this policy has negligible effect on end-to-end FL round time because
local training dominates round duration.

Top-K sparsification is kept optional because it changes the optimization
trajectory, whereas byte-level compression is lossless. When Top-K is enabled,
the omitted residual is accumulated locally:

\begin{equation}
    r_{t+1} = r_t + \Delta\theta_t - TopK(\Delta\theta_t + r_t),
    \label{eq:error_feedback}
\end{equation}

and added back into the next candidate update. This error-feedback mechanism
prevents systematic loss of small but persistent gradients. In the final
evaluation, the communication benchmark focuses on serialization and byte-level
compression because those transformations preserve the model update exactly and
are therefore easier to compare across payload formats.

\section{Causal Root Cause Analysis}
\label{sec:rca_graphs}

Anomaly events are enriched through a causal dependency graph built from
distributed traces. The graph is defined as
$G=(\mathcal{S},\mathcal{E},W)$, where $\mathcal{S}$ is the service set,
$\mathcal{E}$ contains directed calls, and $W$ stores dynamic edge weights. A
trace span pair from service $S_i$ to $S_j$ creates or updates edge
$(S_i,S_j)$. Edge weights use exponential decay:

\begin{equation}
    w_{ij}^{(t+1)} = \lambda w_{ij}^{(t)} + c_{ij}^{(t)},
    \label{eq:edge_decay}
\end{equation}

where $c_{ij}^{(t)}$ is the number of calls observed in interval $t$ and
$\lambda$ discounts stale dependencies.

When anomaly events arrive, the coordinator batches them into an anomalous set
$\mathcal{A}(t)$. A service is considered a stronger root-cause candidate when
it has low anomalous upstream support but high downstream impact. The upstream
anomaly ratio is:

\begin{equation}
    \phi(S_i) =
    \frac{|\mathcal{U}_h(S_i) \cap \mathcal{A}(t)|}
    {|\mathcal{U}_h(S_i)| + \epsilon},
    \label{eq:upstream_ratio}
\end{equation}

where $\mathcal{U}_h(S_i)$ is the upstream neighborhood within $h$ hops. The
impact score combines downstream reachability and edge strength:

\begin{equation}
    I(S_i) = |\mathcal{D}_h(S_i)| +
    \sum_{S_j \in \mathcal{N}^+(S_i)} w_{ij}.
    \label{eq:impact_score}
\end{equation}

The implementation ranks root-cause candidates using graph centrality and impact
evidence, then attaches the top candidates and propagation paths to the alert.
This keeps the RCA output useful for operations: the system does not claim a
single certain cause, but provides a ranked investigation order.

For incidents with multiple anomalous services, the scorer combines structural
centrality with the operational impact score. A personalized PageRank vector
$r$ is computed over the causal graph:

\begin{equation}
    r = (1-\alpha)v + \alpha P^\top r,
    \label{eq:personalized_pagerank}
\end{equation}

where $P$ is the normalized transition matrix and $v$ biases the walk toward
currently anomalous services. The final candidate score is:

\begin{equation}
    Score(S_i) = \alpha_r r_i + \alpha_I \hat{I}(S_i)
    - \alpha_U \phi(S_i),
    \label{eq:rca_final_score}
\end{equation}

where $\hat{I}(S_i)$ is normalized downstream impact and $\phi(S_i)$ penalizes
services that already have anomalous upstream dependencies. This scoring rule
encodes the intuition that root causes tend to be upstream, explanatory, and
high-impact, while propagated anomalies are often downstream symptoms.

\section{Adaptive Thresholding with Q-Learning}
\label{sec:threshold_implementation}

The adaptive thresholding layer is implemented by a per-service
\texttt{ThresholdTuner}. Its state summarizes the current threshold, recent
false-positive rate, recent false-negative rate, and service-level objective
status:

\begin{equation}
    x_t =
    [\tau_t,\text{FPR}_t,\text{FNR}_t,\text{SLO}_t].
    \label{eq:threshold_state}
\end{equation}

The action space is discrete:

\begin{equation}
    \mathcal{U} = \{-0.5,-0.2,0,+0.2,+0.5\},
    \label{eq:threshold_actions}
\end{equation}

with the updated threshold clamped to a minimum of 0.1:

\begin{equation}
    \tau_{t+1} = \max(0.1,\tau_t + u_t).
    \label{eq:threshold_update}
\end{equation}

The Q-table stores $Q(x,u)$ for discretized state buckets. Updates follow the
standard Bellman rule:

\begin{equation}
    Q(x_t,u_t) \leftarrow Q(x_t,u_t) +
    \eta
    \left[
    R_{t+1} + \gamma \max_{u'} Q(x_{t+1},u') - Q(x_t,u_t)
    \right].
    \label{eq:threshold_bellman}
\end{equation}

The reward is asymmetric because missed incidents are more costly than noisy
alerts:

\begin{equation}
    R_t =
    (\alpha_p+\alpha_r)\mathbb{I}(\text{TP})
    - \beta\mathbb{I}(\text{FP})
    - \gamma_{\text{FN}}\mathbb{I}(\text{FN})
    - \delta\mathbb{I}(\text{SLO violation})
    + \delta'\mathbb{I}(\text{SLO compliance}).
    \label{eq:threshold_reward}
\end{equation}

Here, $R_t$ is the scalar feedback assigned after a threshold decision.
$\mathbb{I}(\cdot)$ is an indicator function that equals 1 when the event inside
the parentheses occurs and 0 otherwise. TP, FP, and FN denote true positive,
false positive, and false negative detection outcomes. The coefficients
$\alpha_p$ and $\alpha_r$ reward correct anomaly detection, $\beta$ penalizes
false alarms, $\gamma_{\text{FN}}$ penalizes missed anomalies, $\delta$
penalizes service-level objective violations, and $\delta'$ gives a small
reward for compliant intervals. Table~\ref{tab:reward_coefficients} lists the
values used in the evaluation.

\begin{table}[H]
    \centering
    \small
    \begin{tabular}{|l|c|l|}
        \hline
        \textbf{Coefficient} & \textbf{Value} & \textbf{Rationale} \\
        \hline
        $\alpha_p$ (TP reward) & 1.0 & baseline correct detection reward \\
        \hline
        $\alpha_r$ (recall reward) & 1.0 & equal weight on precision and recall \\
        \hline
        $\beta$ (FP penalty) & 0.5 & false alarms costly but not critical \\
        \hline
        $\gamma_{\text{FN}}$ (FN penalty) & 2.0 & missed incidents more costly than false alarms \\
        \hline
        $\delta$ (SLO violation) & 3.0 & SLO breach has direct business impact \\
        \hline
        $\delta'$ (SLO compliance) & 0.2 & small reward for sustained compliance \\
        \hline
    \end{tabular}
    \caption{Reward coefficients for Q-learning thresholding
             (Eq.~\ref{eq:threshold_reward}). The asymmetry
             $\gamma_{\text{FN}} > \beta$ reflects the operational
             preference that missed anomalies are more harmful than
             false alarms.}
    \label{tab:reward_coefficients}
\end{table}

Reproducing the threshold experiment therefore
requires logging both detection labels or delayed incident feedback and SLO
status; without those signals the Q-table receives only partial feedback.

The evaluation benchmark (10\,020-minute SMD stream, seed~42) shows that the
default tabular reward is not yet stable enough for production use. The
Q-learning policy eventually reduces final alert volume to 28~alerts with
FPR~$=$~0.0155 and F1~$=$~0.5526, but it first passes through a high-alert
exploration phase. This is an important outcome: the adaptive component is
implemented and can improve after feedback, but the reward design must be
strengthened through warm-started Q-tables, constrained exploration, or
DQN-based experience replay.

This finding changes how the thresholding module should be interpreted: the
Q-learning agent is a valuable adaptive control hook that can store per-service
policies and incorporate delayed feedback, but production use should include
warm-start policies from historical incidents or asymmetric reward schedules
where $\gamma_{\text{FN}}$ grows when true anomalies are missed.

\section{Algorithmic Workflow}
\label{sec:method_algorithms}

The workflow has two cooperating loops, summarized in
Algorithms~\ref{alg:edge_agent} and~\ref{alg:coordinator}.

\begin{algorithm}[H]
\caption{Edge Agent Loop}
\label{alg:edge_agent}
\begin{algorithmic}[1]
\Require global model $\theta$, window size $W$, threshold $\tau$
\Loop
    \State $x_t \gets$ \Call{CollectMetrics}{}
    \State $x'_t \gets$ \Call{Normalize}{$x_t$, rolling $\mu$, $\sigma$}
    \State Append $x'_t$ to window buffer
    \If{buffer size $\geq W$}
        \State $X_t \gets$ last $W$ observations
        \State $e_t \gets \frac{1}{WF}\|X_t - \text{Dec}(\text{Enc}(X_t))\|^2$
        \If{$e_t > \tau_t$}
            \State Emit anomaly event $(\text{service}, t, e_t, \tau_t)$
        \EndIf
        \If{labels or SLO feedback available}
            \State $\tau_{t+1} \gets$ \Call{QUpdate}{$\tau_t$, FPR, FNR, SLO}
        \EndIf
    \EndIf
    \If{FL round active}
        \State $\theta_{\text{local}} \gets$ \Call{Train}{$\theta$, local data, $E$ epochs}
        \State $\Delta\theta \gets \theta_{\text{local}} - \theta$
        \State $\Delta\theta \gets$ \Call{ClipAndNoise}{$\Delta\theta$, $C$, $\sigma$}
        \State $\Delta\theta \gets$ \Call{Compress}{$\Delta\theta$}
        \State Submit $\Delta\theta$ to coordinator
    \EndIf
\EndLoop
\end{algorithmic}
\end{algorithm}

\begin{algorithm}[H]
\caption{Coordinator Loop}
\label{alg:coordinator}
\begin{algorithmic}[1]
\Require minimum participants $K_{\min}$, round timeout $T$
\Loop
    \State Open round $t$; record expected model version
    \State Wait for updates until timeout $T$ or $K_{\min}$ received
    \If{$|\text{valid updates}| < K_{\min}$}
        \State Abort round; keep $\theta_t$
    \Else
        \State Validate each update (version, shape, checksum)
        \State $\theta_{t+1} \gets \sum_{k} \frac{n_k}{\sum_j n_j} \theta_t^{(k)}$ \Comment{FedAvg}
        \State Persist $\theta_{t+1}$ to model registry
    \EndIf
    \State Ingest trace spans; update causal graph weights
    \If{anomaly batch window closed}
        \State $\mathcal{A}(t) \gets$ batched anomaly events
        \State Rank root-cause candidates via Eq.~\ref{eq:rca_final_score}
        \State Store enriched alert with graph snapshot
    \EndIf
\EndLoop
\end{algorithmic}
\end{algorithm}

On the edge, the agent collects and normalizes metrics, computes reconstruction
error once $W$ samples are buffered, emits anomaly events when $e_t>\tau_t$, and
updates the threshold only when labels or SLO feedback are available. On the
coordinator, incoming updates are validated (version, shape, checksum) before
FedAvg aggregation. If fewer than $K_{\min}$ valid updates arrive before
timeout, the round is aborted. The coordinator also ingests trace batches and
updates causal graph edge weights.

The methodology presented in this chapter defines the algorithmic core of the
framework. Chapter~\ref{ch:system_design} translates these components into a
deployable system architecture with explicit failure boundaries, persistence
contracts, and configuration surfaces.
